\section{Proofs of the Restated Propagation Results}
\label{app:proofs}

\noindent Every result \emph{MASGym} claims---the metric characterization
(\Cref{thm:metric}), the reduction to ASB (\Cref{prop:reduction}), evaluator integrity
(\Cref{thm:integrity}), judge hijackability (\Cref{prop:judgehijack}), the collusion
advantage (\Cref{thm:collusion}), the forced-coordination bound (\Cref{prop:forced}), and
the concentration and sweep-budget bounds (\Cref{thm:sample,cor:budget})---is stated and
proved \emph{in full in the main paper} (\Cref{sec:theory}), so that no conclusion of this
paper depends on this appendix.

What follows are full derivations of the \emph{classical} spread results that
\Cref{sec:theory-prop} restates in our notation. They are reproduced here for
completeness and are not claimed as novel; the primary statements are Granovetter's
threshold cascades~\cite{granovetter1978threshold}, Watts's mean-field
analysis~\cite{watts2002simple}, Kempe et~al.\ on influence
maximization~\cite{kempe2003maximizing}, and bootstrap
percolation~\cite{chalupa1979bootstrap}. Throughout we work in the independent-cascade
instance of \Cref{as:prop}: when an agent becomes compromised, it independently
compromises each susceptible out-neighbour with probability $\pinf$ when its message is
processed, and compromise is absorbing.

\subsection{Statements}

\begin{proposition}[Chain: bounded reach]\label{prop:chain}
On a directed path $a_1\to\cdots\to a_n$ with a single seed at $a_1$, the expected number
of additionally compromised agents is $\sum_{j=1}^{n-1}\pinf^{\,j}\le \pinf/(1-\pinf)$ for
$\pinf<1$, independent of $n$.
\end{proposition}

\begin{proposition}[Star: the orchestrator is the critical node]\label{prop:star}
On a star with centre $c$, a seed at the centre compromises each leaf independently with
probability $\pinf$, giving expected reach $\pinf(n-1)$ and $\CPR=\pinf$; a seed at a
leaf gives expected reach $\pinf+\pinf^{2}(n-2)$. The centre-seed reach exceeds the
leaf-seed reach by $\Theta(\pinf(1-\pinf)n)$.
\end{proposition}

\begin{proposition}[Tree: a branching phase transition]\label{prop:tree}
On a rooted tree with branching factor $b$ and depth $d$, seeded at the root, the expected
number compromised at level $\ell$ is $(b\pinf)^{\ell}$; the total expected reach is
$\sum_{\ell=1}^{d}(b\pinf)^{\ell}$---bounded in $d$ for $b\pinf<1$ and
$\Theta((b\pinf)^{d})$ for $b\pinf>1$, transitioning at $b\pinf=1$.
\end{proposition}

\begin{proposition}[Mesh: one-round lower bound]\label{prop:mesh}
On the complete graph $K_n$ with $s$ seeds, the expected number compromised after one
round is at least $s+(n-s)\bigl(1-(1-\pinf)^{s}\bigr)$. Consequently, for any fixed
$\pinf>0$ the honest-compromise fraction after one round tends to $1$ as $s$ grows, and
iterating drives $\CPR\to1$ above the percolation threshold.
\end{proposition}

\begin{corollary}[Threshold confinement]\label{cor:threshold}
Under the threshold instance of \Cref{as:prop} with threshold $\thresh$, any agent whose
in-degree in $\Gtop$ is strictly less than $\thresh$ can never be compromised by
propagation. In particular, on a chain (in-degree $1$) with $\thresh\ge2$ the compromised
set never grows beyond the seeds, so $\CPR=0$.
\end{corollary}

\subsection{Derivations}

\begin{proof}[Derivation of \Cref{prop:chain} (chain)]
On the directed path $a_1\to a_2\to\cdots\to a_n$ with the single seed at $a_1$, the agent
$a_{1+j}$ becomes compromised iff every one of the $j$ edges fires. By independence these
events have probability $\pinf^{\,j}$. Summing the expected indicators,
\[
\E\bigl[|\Advset_\infty|-1\bigr]=\sum_{j=1}^{n-1}\pinf^{\,j}
=\pinf\,\frac{1-\pinf^{\,n-1}}{1-\pinf}\le \frac{\pinf}{1-\pinf},
\qquad \pinf<1,
\]
which is bounded independently of $n$.
\end{proof}

\begin{proof}[Derivation of \Cref{prop:star} (star)]
Let the centre be $c$ and the leaves $a_1,\dots,a_{n-1}$.
\emph{Centre seed.} Each leaf is a direct out-neighbour of $c$, so it is compromised
independently with probability $\pinf$; hence $\E[|\Advset_\infty|-1]=\pinf(n-1)$ and
$\CPR=\pinf$.
\emph{Leaf seed.} A seed leaf compromises the centre with probability $\pinf$; conditioned
on that, each of the remaining $n-2$ leaves is compromised with probability $\pinf$. Thus
the expected additional compromises are $\pinf+\pinf^{2}(n-2)$. The centre-seed reach minus
the leaf-seed reach is
\[
\bigl[1+\pinf(n-1)\bigr]-\bigl[1+\pinf+\pinf^{2}(n-2)\bigr]
=\pinf(1-\pinf)(n-2)=\Theta\!\bigl(\pinf(1-\pinf)n\bigr).
\]
This is the star computation that \Cref{thm:collusion} in the main paper turns into the
collusion-advantage result; there we give the proof directly, so the main paper does not
depend on this appendix.
\end{proof}

\begin{proof}[Derivation of \Cref{prop:tree} (tree)]
Seed the root (level $0$). Each compromised node at level $\ell-1$ independently
compromises each of its $b$ children with probability $\pinf$, so the expected number of
compromised nodes at level $\ell$ equals $b\pinf$ times that at level $\ell-1$. With one
seed, $\E[\text{level }\ell]=(b\pinf)^{\ell}$, and the total expected reach is
$\sum_{\ell=1}^{d}(b\pinf)^{\ell}$. For $b\pinf<1$ this is at most $b\pinf/(1-b\pinf)$,
bounded in $d$; for $b\pinf>1$ it is $\Theta\bigl((b\pinf)^{d}\bigr)$; the critical case is
$b\pinf=1$. (This is the mean of a Galton--Watson branching process with offspring mean
$b\pinf$.)
\end{proof}

\begin{proof}[Derivation of \Cref{prop:mesh} (mesh, one round)]
On $K_n$ with seed set $\Advset_0$, $|\Advset_0|=s$, each honest node has all $s$ seeds as
in-neighbours. It survives the round uncompromised iff every one of the $s$ independent
edges fails, with probability $(1-\pinf)^{s}$; hence it is compromised with probability
$1-(1-\pinf)^{s}$. Summing over the $n-s$ honest nodes,
$\E[|\Advset_1|]=s+(n-s)\bigl(1-(1-\pinf)^{s}\bigr)$. For fixed $\pinf>0$,
$1-(1-\pinf)^{s}\to1$ as $s\to\infty$, so the honest-compromise fraction after one round
tends to $1$; iterating the absorbing process only increases the compromised set, and above
the percolation threshold a giant compromised component forms, driving $\CPR\to1$
(cf.~\cite{kempe2003maximizing,watts2002simple}).
\end{proof}

\begin{proof}[Derivation of \Cref{cor:threshold} (threshold confinement)]
In the threshold model an agent is compromised only once at least $\thresh$ of its
in-neighbours are compromised. An agent with fewer than $\thresh$ in-neighbours can never
meet that condition, so it remains honest for the entire (absorbing) process. On a chain
every non-seed agent has in-degree $1<\thresh$, so no propagation occurs.
\end{proof}

\subsection{Why the star is the environment's critical topology}

\Cref{prop:star} and \Cref{thm:collusion} together explain the empirical pattern in
\Cref{sec:defense}: on a star or tree the orchestrator is the single hub through which
every message passes, so one successful injection reaches all workers, whereas on a chain
the reachable set stays bounded by $\pinf/(1-\pinf)$ regardless of length. This is why
\Cref{tab:defense} shows $\ASRs=1.000$ on tree and star but degrades on the chain, and
why \emph{collusion} and \emph{orchestrator compromise} are separate knobs in
\Cref{sec:threat} even though they coincide on the star.